\section{Ratios of marked generating functions}\label{ratios-with-a-fixed-number-of-distinguished-markings}

The comparison theorem controls logarithms, whereas the volume-ratio
problem requires relative errors tending to zero. We obtain that stronger
control here for generating functions with a fixed number of distinguished
markings. The two ratios correspond to adding a cusp and to replacing one
unit of genus by two cusps. Their errors are $O(1/g)$, uniformly in the
parameter. The passage to volumes with a prescribed cusp count will be
made in Section~\ref{extracting-fixed-cusp-counts}.

The mechanism is simple. Positivity of the recursion implies that an
amplitude decreases when one of its degrees is raised~\eqref{eq:TL15}.
The dilaton equation~\eqref{eq:TL12} writes the cusp ratio as a weighted
sum, over degrees \(d\ge1\), of the normalized amplitudes
\(c_{g,p+1}(d,0^p)\). These weights can have either sign; their absolute
moments are uniformly bounded by~\eqref{eq:TL13}. The KdV equation~\eqref{eq:TL14} writes the
genus ratio through \(c_{g,p}(1,0^{p-1})\) and separating terms. Both
ratios are therefore within \(O(1/g)\) of explicit functions of \(Y\) once
we know that \(c_{g,p}(\mathbf d)=1-O(|\mathbf d|^2/g)\). This last estimate,
\eqref{eq:TL23}, follows from the difference of the recursion at adjacent
degrees; the coarse bounds of Section~\ref{uniform-comparison} are used
only to show that the separating terms are small.

\subsection{Shifted amplitudes and normalization}\label{shifted-amplitudes-and-normalization}

Use \(p\) for the number of distinguished markings, and keep \(n\) for the number of cusps. Retain the kernels and shifted amplitudes of Sections~\ref{the-positive-compact-family-of-kernels} and~\ref{exact-amplitudes-and-the-harmonic-model}; write \(U_{g,p}(\mathbf d;Y)=U^Y_{g,p}(\mathbf d)\), and suppress \(Y\) when convenient. Put \(h=2g-2+p\) and \(m=3g-3+p-|\mathbf d|\). Define \begin{equation}
\overline U_{g,p}(\mathbf d)=a^h r^m U_{g,p}(\mathbf d),
 \qquad W_{g,p}=4^g\overline U_{g,p}(\mathbf0),
 \qquad c_{g,p}(\mathbf d)=
 \frac{\overline U_{g,p}(\mathbf d)}{\overline U_{g,p}(\mathbf0)}.
 \label{eq:TL9}
\end{equation} Thus $\overline U_{g,p}(\mathbf d)$ is the normalized amplitude $\mathcal U_{g,p}(\mathbf d;Y)$ of \eqref{eq:H7}; the bar indicates that the markings are now the distinguished ones. The barred amplitudes obey the exact positive recursion \eqref{eq:H9}.

We use two consequences of
Sections~\ref{the-positive-compact-family-of-kernels}--\ref{uniform-comparison}.

\begin{itemize}
\item
  The kernel bounds \eqref{eq:P11}: for each fixed \(k\), \begin{equation}
1=\theta_0\le\theta_1\le\cdots,
   \qquad w_\ell:=\theta_\ell-\theta_{\ell-1}\ge0,
   \quad\theta_{-1}=0,
   \quad\sup_Y\sum_{\ell\ge0}(1+\ell)^k w_\ell<\infty
   \label{eq:TL10}
\end{equation} Also \(A(Y)\) is bounded above and away from
  zero; see Subsection~\ref{compact-normalization-and-uniform-lattice-estimates}.
\item
  Uniform harmonic comparison: for each fixed positive \(p\), \begin{equation}
\log W_{g,p}=
   \log\Gamma(2g-2+p)-(2g-2+p)\log A(Y)+o(g),
   \label{eq:TL11}
\end{equation} uniformly for \(0\le Y<y_0\). The corresponding crude bound
  \(W_{g,p}\le C_p^{g+1}(2g+p)!\) holds uniformly, including bounded
  stable genera. Proposition~\ref{prop:harmonic}, together with the
  central local limit \eqref{eq:H17} for the model, proves
  \eqref{eq:TL11} without any input on volume ratios. Only its \(o(g)\)
  error is needed below, and only for a positive number of markings.
\end{itemize}

\subsection{Dilaton and KdV identities}\label{exact-dilaton-and-genus-exchange-identities}

The shifted weighted dilaton identity is \begin{equation}
h\overline U_{g,p}(\mathbf0)=
 \sum_{d\ge1}b_d(Y)\overline U_{g,p+1}(d,0^p),
 \qquad
 b_d(Y)=\frac{X^{(d)}(Y)r^{d-1}}{a(2d+1)!!}
 =\frac13[z^{d-1}]\bar H_Y(z).
 \label{eq:TL12}
\end{equation} It holds also at \(p=0\), \(g\ge2\). To track the shift, the
background insertion of degree one contributes \(1-a\); moving it and
the higher background insertions to the left leaves the coefficients
\(X^{(d)}/(2d+1)!!\). The factor in \eqref{eq:TL9} then produces exactly
\(r^{d-1}/a\) in \eqref{eq:TL12}. The compact analytic family gives \begin{equation}
\sum_{d\ge1}b_d=A(Y),\qquad
 \sup_Y\sum_{d\ge1}(1+d)^k|b_d(Y)|<\infty
 \label{eq:TL13}
\end{equation} for every fixed \(k\).

The genus ratio is governed instead by the KdV equation. Write $F$ for the Kontsevich--Witten free energy, with subscripts indicating differentiation with respect to the corresponding times. The integrated KdV identity \(F_{01}=F_{00}^2/2+F_{0000}/12\),
differentiated \(p-2\) times in the string variable, gives for \(p\ge2\)
\begin{equation}
\overline U_{g,p}(1,0^{p-1})
 =\frac14\overline U_{g-1,p+2}(\mathbf0)
 +\frac32\sum_{j=0}^g\sum_{s=0}^{p-2}\binom{p-2}s
 \overline U_{j,s+2}(\mathbf0)
 \overline U_{g-j,p-s}(\mathbf0).
 \label{eq:TL14}
\end{equation} The split is ordered. The factor three from the weighted degree-one
insertion turns the original KdV coefficients into \(1/4\) and \(3/2\).
In \eqref{eq:TL14}, every term has the same total powers of \(a\) and \(r\): the
two separating complexities sum to \(h\), and their deficits sum to
\(3g-4+p\). Thus no omitted parameter factor is present.

The key consequence of positivity, used throughout this section, is the
following monotonicity.

\begin{lemma}[Monotonicity in the degrees]\label{lem:monotone}
For every stable \((g,p)\), every degree tuple \(\mathbf d\), every \(i\) and
every \(Y\in[0,y_0]\),
\begin{equation}
\overline U_{g,p}(\mathbf d)\ge
 \overline U_{g,p}(\mathbf d+\mathbf e_i),
 \qquad 0\le c_{g,p}(\mathbf d)\le1.
 \label{eq:TL15}
\end{equation}
Amplitudes outside the dimension range are zero.
\end{lemma}

\begin{proof}
We subtract the recursion at adjacent degrees. With
$\mathbf d'=(d_2,\ldots,d_p)$ fixed, let
$\mathcal B_{\mathbf d'}(t)$ denote the merging and contraction sums
in \eqref{eq:H9}, with the kernel weight removed and $d_1+\ell=t$.
These sums are nonnegative and vanish for
$t>3g-3+p-|\mathbf d'|$ by dimension. If
$m=3g-3+p-|\mathbf d|\ge0$, shifting the kernel index gives
\[
 \begin{aligned}
 \overline U_{g,p}(\mathbf d)
   &=\sum_{\ell=0}^m\theta_\ell
                      \mathcal B_{\mathbf d'}(d_1+\ell),\\
 \overline U_{g,p}(\mathbf d+\mathbf e_1)
   &=\sum_{\ell=1}^m\theta_{\ell-1}
                      \mathcal B_{\mathbf d'}(d_1+\ell).
 \end{aligned}
\]
Their difference is therefore
$\sum_{\ell=0}^m w_\ell\mathcal B_{\mathbf d'}(d_1+\ell)\ge0$, which proves
\eqref{eq:TL15} for \(i=1\), and by symmetry for every \(i\). This includes
the dimension boundary by the zero convention. The difference identity is
used for $h\ge2$; the finite bases satisfy the inequality directly because
\(\theta_1\ge\theta_0=1\).
\end{proof}

From \eqref{eq:TL12}, \eqref{eq:TL13}, and \eqref{eq:TL15}, \begin{equation}
\frac{W_{g,p}}{W_{g,p+1}}\le\frac C{2g-2+p}
 \label{eq:TL16}
\end{equation} in all stable cases used below, with a constant independent of \(p\): the absolute sum in \eqref{eq:TL13} is uniform and the ratios in \eqref{eq:TL15} are at most one. Finite stable bases are included in the constant. Multiplication of \eqref{eq:TL14} by \(4^g\) and use of \eqref{eq:TL15} give \begin{equation}
W_{g-1,p+2}\le W_{g,p},\qquad p\ge2.
 \label{eq:TL17}
\end{equation} The factor \(4^g\) in \eqref{eq:TL9} is essential: it converts the coefficient
\(1/4\) into the comparison constant one.

\subsection{Separating terms}\label{fixed-marking-splitting-estimates}

The separating terms in the KdV identity \eqref{eq:TL14}, and in the
degree difference above, are controlled by the convolution sums \(S_{g,p}\)
and \(T_{g,p}\) below. They are small for two reasons. Splitting off a
component of bounded genus \(j\) removes \(2j+s\) units of complexity,
and the cusp insertions cost a factor \(O(1/g)\) each, by
\eqref{eq:TL16}--\eqref{eq:TL17}; and a split
into two components of comparable genus is exponentially small by the
binomial entropy in \eqref{eq:TL21}.

For fixed \(p\ge2\), define \[
 S_{g,p}=\frac1{W_{g,p}}
 \sum_{j=0}^g\sum_{s=0}^{p-1}\binom{p-1}s
 W_{j,s+1}W_{g-j,p-s},
\] \begin{equation}
T_{g,p}=\frac1{W_{g,p}}
 \sum_{j=0}^g\sum_{s=0}^{p-2}\binom{p-2}s
 W_{j,s+2}W_{g-j,p-s}.
 \label{eq:TL18}
\end{equation} We claim, uniformly in \(Y\), \begin{equation}
S_{g,p}=O_p(g^{-2}),\qquad T_{g,p}=O_p(g^{-1}),
 \label{eq:TL19}
\end{equation} with the positive-genus portion of \(T\) bounded by \(O_p(g^{-2})\).
In particular \(T_{g,2}=O(g^{-2})\) because its genus-zero terms are
unstable.

By symmetry orient the smaller genus as \(j\le g/2\). For
\(1\le j\le\varepsilon g\), insert \(2j\) cusps into the larger-genus
component using \eqref{eq:TL16}, apply \eqref{eq:TL17} \(j\) times, and insert the
remaining \(s\) cusps. This gives \begin{equation}
\frac{W_{g-j,p-s}}{W_{g,p}}
 \le (C_p/g)^{2j+s}.
 \label{eq:TL20}
\end{equation} If \(p-s=1\), instead insert \(2j+1\) cusps before raising the genus,
so every genus-comparison step ends with at least two markings;
afterward insert \(s-1\) cusps. This is possible because \(p\ge2\), and
the same exponent results. In the \(T\) convolution both relevant
marking counts are already at least two.

The crude bound following \eqref{eq:TL11} makes the smaller-component
contribution a fixed polynomial in \(j\) times \[
 (2j)!(C_p/g^2)^j g^{-s}.
\] Choose \(\varepsilon>0\) small enough that successive factorial terms
have a ratio bounded by a fixed number below one. Their positive-genus
sum is \(O_p(g^{-2})\), also with any fixed polynomial in \(j\). For
genus zero, stability gives \(s\ge2\) in \(S\) and \(s\ge1\) in \(T\);
direct use of \eqref{eq:TL16} gives the respective orders.

For \(\varepsilon g\le j\le g/2\), apply \eqref{eq:TL11} to both factors and the
parent. The marking counts belong to a fixed finite set. The powers of
\(A(Y)\) cancel up to bounded exponents, and Stirling's formula gives \begin{equation}
\log\frac{W_{j,s+c}W_{g-j,p-s}}{W_{g,p}}
 \le-2g\mathsf H(j/g)+o(g),\qquad c\in\{1,2\},
 \label{eq:TL21}
\end{equation} where \(\mathsf H(t)=-t\log t-(1-t)\log(1-t)\). Its positive minimum
on \([\varepsilon,1/2]\) makes the macroscopic sum exponentially small.
This proves \eqref{eq:TL19}. Every use of \eqref{eq:TL11} here has a positive fixed
marking count.

\subsection{Cusp and genus ratios}\label{uniform-fixed-marking-ratios}

For \(p\ge2\), the adjacent-degree difference obtained from \eqref{eq:H9} has a
merging prefactor \(O_p(1/g)\) by \eqref{eq:TL16}. Its nonseparating prefactor is
bounded by \begin{equation}
\frac{W_{g-1,p+1}}{W_{g,p}}
 \le\frac Cg\frac{W_{g-1,p+2}}{W_{g,p}}=O_p(g^{-1})
 \label{eq:TL22}
\end{equation} using \eqref{eq:TL16}--\eqref{eq:TL17}. To see all
three contributions, put $D=|\mathbf d|$,
$\Omega_0=\sum_{\ell\ge0}w_\ell$ and
$\Omega_1=\sum_{\ell\ge0}\ell w_\ell$. Bound each child amplitude
by its zero-degree value in the difference identity above. Since there
are at most $d_1+\ell+1$ new-degree pairs, it yields
\[
 \begin{aligned}
 c_{g,p}(\mathbf d)-c_{g,p}(\mathbf d+\mathbf e_1)
 \le{}&\Omega_0(2D+p-1)\frac{W_{g,p-1}}{W_{g,p}}\\
 &+\bigl((d_1+1)\Omega_0+\Omega_1\bigr)
   \left(2\frac{W_{g-1,p+1}}{W_{g,p}}+\frac12S_{g,p}\right).
 \end{aligned}
\]
The factor $2$ in the nonseparating term includes the factor $4$
from the genus normalization of $W$. The two kernel moments are
uniformly bounded by \eqref{eq:TL10}. Equations~\eqref{eq:TL16},
\eqref{eq:TL19} and~\eqref{eq:TL22} now give \[
 0\le c_{g,p}(\mathbf d)-c_{g,p}(\mathbf d+\mathbf e_i)
 \le\frac{C_p(|\mathbf d|+1)}g.
\] Summing degree increments proves \begin{equation}
0\le1-c_{g,p}(\mathbf d)
 \le C_p\frac{|\mathbf d|(|\mathbf d|+1)}g,
 \qquad p\ge2.
 \label{eq:TL23}
\end{equation} When the right side exceeds one the bound is trivial. It also
bounds the terms outside the dimension range when they are set to zero.

Substitute \eqref{eq:TL23} into \eqref{eq:TL12}. The errors sum absolutely by \eqref{eq:TL13},
giving \begin{equation}
(2g-2+p)\frac{W_{g,p}}{W_{g,p+1}}
 =A(Y)+O_p(g^{-1}),\qquad p\ge1.
 \label{eq:TL24}
\end{equation} Since \(A\) is uniformly positive, the error is relative
\(O_p(1/g)\).

To recover the one-marking profile without circularity, first use \eqref{eq:TL16}
twice, \eqref{eq:TL17} at parent marking count two, and the now two-sided ratio
\eqref{eq:TL24} at marking count one: \begin{equation}
\frac{W_{g-1,2}}{W_{g,1}}
 \le\frac C{g^2}\frac{W_{g-1,4}}{W_{g,2}}
 \frac{W_{g,2}}{W_{g,1}}=O(g^{-1}).
 \label{eq:TL25}
\end{equation} The separating terms at one marking have two positive genera.
Equation \eqref{eq:TL16}, the inequality \(1/[j(g-j)]\le C/g\), and
\(T_{g,2}=O(g^{-2})\) give \begin{equation}
\frac1{W_{g,1}}\sum_{j=1}^{g-1}W_{j,1}W_{g-j,1}
 \le\frac Cg\frac{W_{g,2}}{W_{g,1}}T_{g,2}=O(g^{-2}).
 \label{eq:TL26}
\end{equation} The same degree-difference argument now proves \eqref{eq:TL23} for \(p=1\).
Thus \eqref{eq:TL24} holds also for \(p=0\).

Finally \eqref{eq:TL14}, \eqref{eq:TL19}, and \eqref{eq:TL23} yield \begin{equation}
\frac{W_{g-1,p+2}}{W_{g,p}}=1+O_p(g^{-1}),\qquad p\ge2.
 \label{eq:TL27}
\end{equation} For \(p=0,1\), add two markings in numerator and denominator and use
\eqref{eq:TL24}; the two cusp products have the same complexity factors and
cancel. Hence \eqref{eq:TL27} holds for every fixed \(p\ge0\).

We now return to volumes. Recall the cusp generating function
\(\Phi_g\) of \eqref{eq:intro-Phi} and the string-shift identity
\eqref{eq:H8}: \(U_{g,p}(0^p;Y)=\Phi_g^{(p)}(x)\) with \(x=X(Y)\), whenever
the amplitude is stable. By \eqref{eq:TL9},
\(W_{g,p}=4^ga^hr^m\Phi_g^{(p)}(x)\) with \(m=3g-3+p\). Adding a marking
raises both \(h\) and \(m\) by one, so
\[
 \frac{x\Phi_g^{(p+1)}(x)}{\Phi_g^{(p)}(x)}
 =\frac{x}{ar}\,\frac{W_{g,p+1}}{W_{g,p}}
 =\frac{hx}{A\,a\,r}\bigl(1+O_p(g^{-1})\bigr)
\]
by \eqref{eq:TL24}. Replacing \((g,p)\) by \((g-1,p+2)\) leaves \(h\)
unchanged and lowers \(m\) by one, so
\(\Phi_{g-1}^{(p+2)}/\Phi_g^{(p)}=4r\,W_{g-1,p+2}/W_{g,p}\), and
\eqref{eq:TL27} applies. This proves, for every fixed positive \(p\),
\begin{equation}
x\frac{(\Phi_g^{(p)})'(x)}{\Phi_g^{(p)}(x)}
 =h\Psi(x)+O_p(\Psi(x)),\qquad h=2g-2+p,
 \label{eq:TL30}
\end{equation} \begin{equation}
\frac{\Phi_{g-1}^{(p+2)}(x)}{\Phi_g^{(p)}(x)}
 =4r(Y)(1+O_p(g^{-1})),
 \qquad
 \Psi(x)=\frac{x}{A(Y)a(Y)r(Y)}.
 \label{eq:TL31}
\end{equation} The same genus ratio holds at \(p=0\). The estimates are uniform
throughout \(0\le Y<y_0\); at \(x=0\), \eqref{eq:TL30} is interpreted by
continuity.

The scalar is exactly the stationary parameter: \begin{equation}
H_Y(r)=\frac{3I(Y)}{2\Delta^{3/2}},\qquad
 A(Y)a(Y)r(Y)=\frac{I(Y)}{\sqrt\Delta},\qquad
 \Psi(x)=\lambda(Y)/2.
 \label{eq:TL32}
\end{equation} The first equality follows by substituting \(v=T-\Delta s^2\) in
\(I\), and the other two follow by direct substitution. We summarize the
output of this section used in the extraction arguments.

\begin{proposition}[Marked cusp and genus ratios]\label{prop:marked-ratios}
For every fixed \(p\ge1\), uniformly for \(0\le Y<y_0\) and \(x=X(Y)\), the
cusp ratio of \(\Phi_g^{(p)}\) satisfies \eqref{eq:TL30} and the genus ratio
satisfies \eqref{eq:TL31}, with \(\Psi(x)=\lambda(Y)/2\). The genus ratio
holds also at \(p=0\).
\end{proposition}
